\chapter{Autodeclaração do Uso de Inteligência Artificial}
\label{cap:autodeclaracao-ia}

Em atenção às boas práticas de integridade acadêmica e à transparência no processo de produção científica, os autores declaram que ferramentas de Inteligência Artificial (IA) generativa, baseadas em grandes modelos de linguagem, foram utilizadas como apoio durante a elaboração deste trabalho. O uso dessas ferramentas ocorreu de forma complementar e supervisionada, sem substituir a autoria intelectual, a análise crítica e as decisões dos autores.

As ferramentas de IA foram empregadas nas seguintes atividades:

\begin{itemize}
    \renewcommand\labelitemi{-}
    \item \textbf{Correções gramaticais e revisão textual:} identificação de erros de ortografia, concordância, pontuação e coesão, bem como sugestões de reescrita para tornar o texto mais claro e adequado à linguagem acadêmica.
    \item \textbf{Geração e discussão de ideias:} apoio na organização de argumentos, na estruturação de seções e na exploração de alternativas para a redação de trechos do trabalho.
    \item \textbf{Pesquisa de autores e referências:} auxílio na localização de autores, artigos e documentos relacionados ao tema, servindo como ponto de partida para a busca bibliográfica.
    \item \textbf{Dúvidas sobre \textit{Overleaf} e \LaTeX:} esclarecimento sobre comandos, formatação, compilação, organização de referências, listas e demais recursos necessários à editoração do documento conforme as normas da instituição.
\end{itemize}

Todo conteúdo sugerido pelas ferramentas foi analisado, ajustado e validado pelos autores antes de ser incorporado ao texto. Em especial, as referências e os dados indicados com o auxílio da IA foram conferidos diretamente nas fontes originais, uma vez que tais ferramentas podem produzir informações imprecisas ou inexistentes. Sugestões que não puderam ser confirmadas em fontes confiáveis foram descartadas ou reformuladas.

É importante ressaltar que a IA não foi utilizada para gerar os dados, os resultados ou as conclusões apresentados neste trabalho. A concepção do projeto, o levantamento de requisitos junto aos profissionais da área veterinária, o desenvolvimento da aplicação RoadVet e a análise dos resultados são de inteira responsabilidade dos autores, que respondem integralmente pelo conteúdo deste documento.
